\documentclass[10pt,a4paper]{article}
\usepackage[top=1.5cm, bottom=1.5cm, left=1.6cm, right=1.6cm, headheight=14pt, footskip=18pt]{geometry}
\usepackage[utf8]{inputenc}
\usepackage{amsmath,amsfonts,amssymb}
\usepackage{graphicx}
\usepackage{tikz}
\usepackage{circuitikz}
\usetikzlibrary{calc,patterns}
\usepackage[most]{tcolorbox}
\usepackage{enumitem}
\usepackage{fancyhdr}
\usepackage{booktabs}
\usepackage{tabularx}
\usepackage{colortbl}
\usepackage{hyperref}

% --- Palet Warna Resmi UNIB ---
\definecolor{unibblue}{RGB}{0, 32, 96}
\definecolor{unibgold}{RGB}{197, 160, 89}
\definecolor{unibgreen}{RGB}{34, 112, 60}
\definecolor{boxbg}{RGB}{248, 250, 253}
\definecolor{linegray}{RGB}{210, 215, 225}

% --- Pengaturan Header & Footer ---
\pagestyle{fancy}
\fancyhf{}
\lhead{\footnotesize\textbf{\color{unibblue}Universitas Bengkulu} \textbar\ Program Studi S1 Teknik Elektro}
\rhead{\footnotesize\color{gray}Persamaan Diferensial (Genap 2026)}
\lfoot{\footnotesize\color{gray}Lembar Kerja Mahasiswa (LKM) Berbasis OBE}
\cfoot{\footnotesize\color{gray}Hal. \thepage\ dari 2}
\rfoot{\footnotesize\color{unibblue}\textbf{Minggu 7: Invers Transformasi Laplace \& Pe...}}
\renewcommand{\headrulewidth}{0.6pt}
\renewcommand{\footrulewidth}{0.4pt}

% --- Custom Box untuk Lembar Jawab ---
\newtcolorbox{answerbox}[1]{%
  colback=white,
  colframe=unibblue!70!black,
  coltitle=white,
  fonttitle=\bfseries\scriptsize,
  colbacktitle=unibblue,
  title={#1},
  arc=2pt,
  left=5pt,right=5pt,top=3pt,bottom=3pt,
  boxrule=0.6pt
}

\begin{document}

% ==================== HALAMAN 1 ====================
\noindent\begin{minipage}[c]{0.68\textwidth}%
  {\large\textbf{\color{unibblue}LEMBAR KERJA MAHASISWA (LKM)}}\\[1mm]
  {\normalsize\textbf{Mata Kuliah: Persamaan Diferensial (2 SKS)}}\\[0.5mm]
  {\small\textbf{Topik Minggu 7:} Invers Transformasi Laplace \& Pemutus Tenaga (TRV)}%
\end{minipage}%
\hfill%
\begin{minipage}[c]{0.30\textwidth}%
  \begin{tcolorbox}[colback=unibblue!5,colframe=unibblue,boxrule=0.6pt,arc=2pt,left=4pt,right=4pt,top=2pt,bottom=2pt]
    \scriptsize
    \textbf{Nama:} \dotfill\\[1.2mm]
    \textbf{NPM:} \dotfill\\[1.2mm]
    \textbf{Kelas/Klp:} \dotfill\\[1.2mm]
    \textbf{Nilai Final:} \hfill \textbf{/ 100}
  \end{tcolorbox}%
\end{minipage}\par

\vspace{1.5mm}

% Kotak Sub-CPMK & Pedoman
\begin{tcolorbox}[colback=boxbg,colframe=unibgold!90!black,boxrule=0.6pt,arc=2pt,left=5pt,right=5pt,top=3pt,bottom=3pt,title=\textbf{\scriptsize Capaian Pembelajaran (Sub-CPMK 4) \& Pedoman Evaluasi OBE}]
  \scriptsize
  \textbf{Sub-CPMK 4:} Mahasiswa mampu merekonstruksi sinyal waktu melalui Invers Transformasi Laplace (pecahan parsial \& integral konvolusi) serta menganalisis fenomena Transient Recovery Voltage (TRV) pada pemutus tenaga PMT.\\
  \textbf{Alokasi Waktu:} 50--60 Menit \quad\textbar\quad \textbf{Model:} Kolaboratif / Mandiri Terstruktur \quad\textbar\quad \textbf{Bahasa Komputasi:} Julia 1.12+
\end{tcolorbox}

\vspace{1mm}

% Soal C1
\noindent\textbf{1. (C1 - Mengingat \textbar\ Bobot: 10 Poin)}\par\vspace{0.5mm}
{\footnotesize Tuliskan tiga kasus dekomposisi pecahan parsial (*partial fraction expansion*) untuk fungsi rasional $F(s) = P(s)/Q(s)$ (akar riil berbeda, akar riil berulang, dan pasangan akar kompleks konjugat).}
\begin{answerbox}{Ruang Pengerjaan C1 (Kaidah Dekomposisi Pecahan Parsial)}
  \vspace{2.0cm}
\end{answerbox}

\vspace{1mm}

% Soal C2
\noindent\textbf{2. (C2 - Memahami \textbar\ Bobot: 15 Poin)}\par\vspace{0.5mm}
{\footnotesize Jelaskan definisi fisis *Transient Recovery Voltage* (TRV) yang muncul melintasi kontak Pemutus Tenaga (PMT / *Circuit Breaker*) sesaat setelah pemutusan arus hubung singkat menurut standar \textbf{IEC 62271-100}.}
\begin{answerbox}{Ruang Pengerjaan C2 (Intuisi Fisika TRV \& Kegagalan Dielektrik PMT)}
  \vspace{2.1cm}
\end{answerbox}

\vspace{1mm}

% Soal C3
\noindent\textbf{3. (C3 - Menerapkan \textbar\ Bobot: 20 Poin)}\par\vspace{0.5mm}
\noindent\begin{minipage}[t]{0.60\textwidth}%
  {\footnotesize Model domain-$s$ dari tegangan pemulihan transien PMT diberikan oleh $V_{\text{TRV}}(s) = V_m \left( \frac{1}{s} - \frac{s}{s^2 + \omega_0^2} \right)$. Gunakan invers transformasi Laplace untuk menurunkan fungsi waktu analitis $v_{\text{TRV}}(t)$, dan tentukan nilai puncak maksimumnya!}%
\end{minipage}%
\hfill%
\begin{minipage}[t]{0.37\textwidth}%
  \centering
  \resizebox{0.95\linewidth}{!}{%
  \begin{circuitikz}[american, scale=0.65, transform shape]
    \draw (0,0) to[sV, l=$v_s(t)$] (0,1.8)
          to[L, l=$L_{\text{trafo}}$] (1.6,1.8)
          to[switch, l={PMT Buka}] (3.0,1.8)
          to[short] (3.0,0)
          to[short] (0,0);
    \draw (1.6,1.8) to[C, l=$C_{\text{bus}}$, v=$v_{\text{TRV}}$] (1.6,0);
  \end{circuitikz}%
  }%
\end{minipage}\par\vspace{1mm}
\begin{answerbox}{Ruang Pengerjaan C3 (Invers Laplace \& Penentuan Nilai Puncak TRV)}
  \vspace{2.8cm}
\end{answerbox}

\newpage
% ==================== HALAMAN 2 ====================

% Soal C4
\noindent\textbf{4. (C4 - Menganalisis \textbar\ Bobot: 20 Poin)}\par\vspace{0.5mm}
{\footnotesize Analisis laju kenaikan tegangan pemulihan (\textit{Rate of Rise of Recovery Voltage} / RRRV) yang didefinisikan sebagai $\text{RRRV} = \left. \frac{dv_{\text{TRV}}}{dt} \right|_{\text{maks}}$. Jelaskan mengapa jika RRRV melebihi kekuatan dielektrik celah busur gas $\text{SF}_6$, pemutus tenaga akan mengalami *re-strike* (kegagalan interupsi).}
\begin{answerbox}{Ruang Pengerjaan C4 (Analisis Kritis Parameter RRRV Standar IEC 62271)}
  \vspace{3.8cm}
\end{answerbox}

\vspace{1.5mm}

% Soal C5
\noindent\textbf{5. (C5 - Mengevaluasi \textbar\ Bobot: 20 Poin)}\par\vspace{0.5mm}
{\footnotesize Selesaikan PDB sistem kendali eksitasi dengan nilai awal tidak nol: $y'' + 4y' + 13y = 0$, $y(0) = 5$, $y'(0) = -10$. Evaluasi apakah respons waktu mengalami osilasi teredam, dan tentukan frekuensi osilasi teredamnya $\omega_d$.}
\begin{answerbox}{Ruang Pengerjaan C5 (Penyelesaian Lengkap IVP Domain-s \& Evaluasi Redaman)}
  \vspace{3.8cm}
\end{answerbox}

\vspace{1.5mm}

% Soal C6
\noindent\textbf{6. (C6 - Merancang \& Komputasi Julia \textbar\ Bobot: 15 Poin)}\par\vspace{0.5mm}
{\footnotesize Rancang skrip numerik berbasis \textbf{Julia} untuk menyelesaikan PDB TRV orde 2 non-homogen dengan paket \texttt{Differential\-Equations.jl} serta memplot kurva $v_{\text{TRV}}(t)$ bersama garis batas dielektrik PMT.}
\begin{answerbox}{Ruang Pengerjaan C6 (Skrip Simulasi Numerik TRV PMT Gardu Induk Julia)}
  \vspace{3.8cm}
\end{answerbox}

\vspace{1.5mm}

% Tabel Rekapitulasi Skor OBE & Tanda Tangan
\noindent\begin{minipage}[b]{0.65\textwidth}%
  \centering
  \scriptsize
  \begin{tabular}{|c|c|c|c|c|c|c|}
    \hline
    \rowcolor{unibblue}
    \textbf{\color{white}C1} & \textbf{\color{white}C2} & \textbf{\color{white}C3} & \textbf{\color{white}C4} & \textbf{\color{white}C5} & \textbf{\color{white}C6} & \textbf{\color{white}Total} \\
    \rowcolor{unibblue!10}
    10 Pts & 15 Pts & 20 Pts & 20 Pts & 20 Pts & 15 Pts & 100 Pts \\
    \hline
    & & & & & & \\[2.5mm]
    \hline
  \end{tabular}%
\end{minipage}%
\hfill%
\begin{minipage}[b]{0.32\textwidth}%
  \centering
  \scriptsize
  \textbf{Verifikasi Dosen / Asisten:}\\[5mm]
  (\dotfill)
\end{minipage}\par

\end{document}
